% GENERATED FILE. Edit canonical lessons, then run npm run generate:books.
% canonical-source-hash: fnv1a64:56fef0d4c013830e
% canonical-lessons: JA-C147-arau, JA-C147-migaku, JA-C147-abiru, JA-C147-yu, JA-C147-sekken

\chapter{Washing Up with Soap and Hot Water}
\label{ch:ja-washing-up-with-soap-and-hot-water}
\hlchaptermodality{147}

\begin{chapteropening}
\textbf{By the end of this chapter:} \emph{I can say how I wash in the morning.}

Complete the last lesson of chapter 147: I can say how I wash in the morning.
\end{chapteropening}

\section[\texorpdfstring{\ja{あらう} (arau)}{あらう (arau)}]{\ja{あらう} (arau) — to wash}

\label{lesson:JA-C147-arau}



\begin{quote}
Before the new one: say the Japanese for to wake someone, then the Japanese for sleep.
\end{quote}

\subsection*{You'll want to know: \ja{あらう}}
\textbf{\ja{あらう}} — \emph{arau} — \textquotedblleft{}to wash\textquotedblright{}.

\begin{quote}
\textbf{\ja{まいあさ} \ja{かおを} \ja{あらいます。}} — \emph{maiasa kao o araimasu} — \textquotedblleft{}I wash my face every morning.\textquotedblright{}
\end{quote}

Polite \textbf{\ja{あらいます}}, \emph{araimasu}; a godan verb.

\begin{cousinweb}
\textbf{\ja{あらう}} (\emph{arau}) is a native Japanese word (wago), inherited rather than borrowed from Chinese; its kanji was attached later as a kun'yomi, and its older origin is not settled (JMdict, KANJIDIC2).
\end{cousinweb}

\begin{grammarlens}[title={What you've built}]
One more word for everyday life at A2.
\end{grammarlens}

\subsection*{Your turn}
\begin{itemize}
\raggedright
  \item \emph{Say it:} \emph{arau}
  \item \emph{Say it:} \emph{arau}, once more
  \item \emph{Say it:} \emph{maiasa kao o araimasu} — I wash my face every morning.
  \item \emph{Recall:} say the Japanese for to wake someone, then the Japanese for sleep, then say \emph{arau} again
\end{itemize}

\begin{tcolorbox}[breakable,colback=teal!4,colframe=teal!35!black,title={Before you move on}]
What does \ja{あらう} mean? (\textquotedblleft{}To wash\textquotedblright{}.) Say it once more.
\end{tcolorbox}
\section[\texorpdfstring{\ja{みがく} (migaku)}{みがく (migaku)}]{\ja{みがく} (migaku) — to polish; to brush (teeth)}

\label{lesson:JA-C147-migaku}



\begin{quote}
Before the new one: say the Japanese for sleep, then the Japanese for to wash.
\end{quote}

\subsection*{You'll want to know: \ja{みがく}}
\textbf{\ja{みがく}} — \emph{migaku} — \textquotedblleft{}to polish; to brush (teeth)\textquotedblright{}.

\begin{quote}
\textbf{\ja{まいにち} \ja{はを} \ja{みがきます。}} — \emph{mainichi ha o migakimasu} — \textquotedblleft{}I brush my teeth every day.\textquotedblright{}
\end{quote}

Polite \textbf{\ja{みがきます}}, \emph{migakimasu}; a godan verb.

\begin{cousinweb}
\textbf{\ja{みがく}} (\emph{migaku}) is a native Japanese word (wago), inherited rather than borrowed from Chinese; its kanji was attached later as a kun'yomi, and its older origin is not settled (JMdict, KANJIDIC2).
\end{cousinweb}

\begin{grammarlens}[title={What you've built}]
One more word for everyday life at A2.
\end{grammarlens}

\subsection*{Your turn}
\begin{itemize}
\raggedright
  \item \emph{Say it:} \emph{migaku}
  \item \emph{Say it:} \emph{migaku}, once more
  \item \emph{Say it:} \emph{mainichi ha o migakimasu} — I brush my teeth every day.
  \item \emph{Recall:} say the Japanese for sleep, then the Japanese for to wash, then say \emph{migaku} again
\end{itemize}

\begin{tcolorbox}[breakable,colback=teal!4,colframe=teal!35!black,title={Before you move on}]
What does \ja{みがく} mean? (\textquotedblleft{}To polish; to brush (teeth)\textquotedblright{}.) Say it once more.
\end{tcolorbox}
\section[\texorpdfstring{\ja{あびる} (abiru)}{あびる (abiru)}]{\ja{あびる} (abiru) — to pour over oneself; to take (a shower)}

\label{lesson:JA-C147-abiru}



\begin{quote}
Before the new one: say the Japanese for to wash, then the Japanese for to polish.
\end{quote}

\subsection*{You'll want to know: \ja{あびる}}
\textbf{\ja{あびる}} — \emph{abiru} — \textquotedblleft{}to pour over oneself; to take (a shower)\textquotedblright{}.

\begin{quote}
\textbf{\ja{まいあさ} \ja{みずを} \ja{あびます。}} — \emph{maiasa mizu o abimasu} — \textquotedblleft{}Every morning I pour cold water over myself.\textquotedblright{}
\end{quote}

Polite \textbf{\ja{あびます}}, \emph{abimasu}; an ichidan verb.

\begin{cousinweb}
\textbf{\ja{あびる}} (\emph{abiru}) is a native Japanese word (wago), inherited rather than borrowed from Chinese; its kanji was attached later as a kun'yomi, and its older origin is not settled (JMdict, KANJIDIC2).
\end{cousinweb}

\begin{grammarlens}[title={What you've built}]
One more word for everyday life at A2.
\end{grammarlens}

\subsection*{Your turn}
\begin{itemize}
\raggedright
  \item \emph{Say it:} \emph{abiru}
  \item \emph{Say it:} \emph{abiru}, once more
  \item \emph{Say it:} \emph{maiasa mizu o abimasu} — Every morning I pour cold water over myself.
  \item \emph{Recall:} say the Japanese for to wash, then the Japanese for to polish, then say \emph{abiru} again
\end{itemize}

\begin{tcolorbox}[breakable,colback=teal!4,colframe=teal!35!black,title={Before you move on}]
What does \ja{あびる} mean? (\textquotedblleft{}To pour over oneself; to take (a shower)\textquotedblright{}.) Say it once more.
\end{tcolorbox}
\section[\texorpdfstring{\ja{ゆ} (yu)}{ゆ (yu)}]{\ja{ゆ} (yu) — hot water; a hot bath}

\label{lesson:JA-C147-yu}



\begin{quote}
Before the new one: say the Japanese for to polish, then the Japanese for to pour over oneself.
\end{quote}

\subsection*{You'll want to know: \ja{ゆ}}
\textbf{\ja{ゆ}} — \emph{yu} — \textquotedblleft{}hot water; a hot bath\textquotedblright{}.

\begin{quote}
\textbf{\ja{おふろの} \ja{ゆは} \ja{あついです。}} — \emph{ofuro no yu wa atsui desu} — \textquotedblleft{}The bath water is hot.\textquotedblright{}
\end{quote}

\begin{cousinweb}
\textbf{\ja{ゆ}} (\emph{yu}) is a native Japanese word (wago), inherited rather than borrowed from Chinese; its kanji was attached later as a kun'yomi, and its older origin is not settled (JMdict, KANJIDIC2).
\end{cousinweb}

\begin{grammarlens}[title={What you've built}]
One more word for everyday life at A2.
\end{grammarlens}

\subsection*{Your turn}
\begin{itemize}
\raggedright
  \item \emph{Say it:} \emph{yu}
  \item \emph{Say it:} \emph{yu}, once more
  \item \emph{Say it:} \emph{ofuro no yu wa atsui desu} — The bath water is hot.
  \item \emph{Recall:} say the Japanese for to polish, then the Japanese for to pour over oneself, then say \emph{yu} again
\end{itemize}

\begin{tcolorbox}[breakable,colback=teal!4,colframe=teal!35!black,title={Before you move on}]
What does \ja{ゆ} mean? (\textquotedblleft{}Hot water; a hot bath\textquotedblright{}.) Say it once more.
\end{tcolorbox}
\section[\texorpdfstring{\ja{せっけん} (sekken)}{せっけん (sekken)}]{\ja{せっけん} (sekken) — soap}

\label{lesson:JA-C147-sekken}



\begin{quote}
Before the new one: say the Japanese for to pour over oneself, then the Japanese for hot water.
\end{quote}

\subsection*{You'll want to know: \ja{せっけん}}
\textbf{\ja{せっけん}} — \emph{sekken} — \textquotedblleft{}soap\textquotedblright{}.

\begin{quote}
\textbf{\ja{せっけんで} \ja{てを} \ja{あらいます。}} — \emph{sekken de te o araimasu} — \textquotedblleft{}I wash my hands with soap.\textquotedblright{}
\end{quote}

\begin{cousinweb}
\textbf{\ja{せっけん}} (\emph{sekken}) is Sino-Japanese (kango): two Chinese characters read with their on'yomi, readings borrowed from Chinese. They are \emph{seki}, \textquotedblleft{}stone\textquotedblright{} (Middle Chinese \emph{dzyek}) and \emph{ken}, \textquotedblleft{}alkaline salt\textquotedblright{} (\emph{keamX}) (KANJIDIC2; Middle Chinese in Baxter's transcription).
\end{cousinweb}

\begin{grammarlens}[title={What you've built}]
One more word for everyday life at A2.
\end{grammarlens}

\subsection*{Your turn}
\begin{itemize}
\raggedright
  \item \emph{Say it:} \emph{sekken}
  \item \emph{Say it:} \emph{sekken}, once more
  \item \emph{Say it:} \emph{sekken de te o araimasu} — I wash my hands with soap.
  \item \emph{Recall:} say the Japanese for to pour over oneself, then the Japanese for hot water, then say \emph{sekken} again
\end{itemize}

\begin{tcolorbox}[breakable,colback=teal!4,colframe=teal!35!black,title={Before you move on}]
What does \ja{せっけん} mean? (\textquotedblleft{}Soap\textquotedblright{}.) Say it once more.
\end{tcolorbox}
